% !TeX root = ../papers/nc-p3.tex
\section{正則性座標 $x=0$ の全点とその添字保持分類}
\label{sec:rank-fifteen-x-zero}
\begin{theorem}\label{thm:rank-fifteen-x-zero}
標数零の階数15のPauli曲面で、$x=0$ の全点は標準Koszul左右AS正則で、
各始点で四中心四次式の多項式環上自由階数256となる。
この軌跡の添字保持同型類は $j\in\mathbb G_m\setminus\{1,-1\}$ で重複なく表され、
\[
 p=\frac{3\sqrt3}{4}(j-j^{-1}),\qquad
 q=\frac{j+j^{-1}}2,\qquad
 y=-\frac{(j-j^{-1})^2}{16}
\]
である。全ての最小周期は2である。$y$ を固定したとき $y=1/4$ で二類、
それ以外の $y\ne0$ で四類となる。
\end{theorem}
\begin{proof}
$x=0$ では $U=-3$、$H^6=256y$ と取れる。
3の上で $y$ が単元なら定理\ref{thm:rank-fifteen-residue-line}、
極を持てば定理\ref{thm:rank-fifteen-pole-cone}が適用される。
残る場合は $\nu(H)>0$ である。$i^2=-1,d^2=3$ とし、$K=id$ を取る。
$D=\mathbb Z_{(3)}[i,d]$ の $d$ 上の剰余体は $\mathbb F_9$ である。
$D[[H]]$ 上の四つの平方根
\[
 P,P'=\sqrt{2(K+1+H)},\sqrt{2(K+1-H)},\qquad
 Q,Q'=\sqrt{2(K-1+H)},\sqrt{2(K-1-H)}
\]
の定数項をそれぞれ $d+i,d+i,-(1+K),-(1+K)$ とする。
$(d+i)^2=2(K+1)$、$(1+K)^2=2(K-1)$ なので全て単根である。
$n=\sqrt{(-4-\sqrt{16-4H^2})/2}$ は内側の定数項が4、外側が $-2i$ の根とする。
これで $n^4+4n^2+H^2=0$ が成立する。

\path{computations/verify_rank_fifteen_x_zero.py} は $1,i,d,id$ の四座標を有理係数で持つ。
定数項が単元である条件は、全座標が3整で、初めの二座標の少なくとも一つが
3で非零になることである。$H$ で割る際は四座標とも定数項が
有理数として零であることを確認する。
前節と共通の経路・中心・商の行列を精度25から計算した。
長さ4,5の経路での精度損失はそれぞれ2、四次中心方程式では零となる。
文字零の後二自由列と文字一の両自由列を選び、四中心商を各始点で長さ13まで
計算すると、各長さの最大損失は前節と同列で最後の精度は1である。
従って中心の持上げと商次元の上界から全次数の正則性・自由性が従う。
より一般に $\nu(K-id)\ge25\nu(H)>0$ の窓にも同じ計算が適用できる。
実際 $K=id+H^{25}T$ としても四つの平方根と四次根の選んだ解は
精度25まで一致し、行操作に必要な全ての定数項が保たれる。

代数的な $y$ には数体の3の上の付値を取り、超越的な $y$ には
3進Gauss付値を取れば、以上の三場合で全点を尽くす。
商の方程式より $q^2-4p^2/27=1$ である。
$j=q+2p/(3\sqrt3)$ は単元で、その逆元は $q-2p/(3\sqrt3)$ である。
$y\ne0$ は $j\ne\pm1$ に等しい。定理\ref{thm:rank-fifteen-formal-boundary}の
四例外は全て $x\ne0$ なので、ここでは全点が通常表示を持つ。
$(x,p,q)$ は完全な添字保持同型不変量で、一位置移動は $j\mapsto-j$、
二位置移動は同型になる。$p$ は常に非零で、$q=0$ は $y=1/4$ に限られるから
最小周期と各ファイバーの個数も従う。
\end{proof}

\begin{corollary}
剰余標数3で $\alpha=\nu(3)$, $\beta=\nu(y)>0$,
$\delta=\nu(x^3/y)$ とする。$\delta>3\alpha$ かつ
$2\delta\ge3\alpha+25\beta$ ならば、通常な階数15窓は同じ正則性と自由性を持つ。
\end{corollary}
\begin{proof}
$x=0$ は定理に含まれる。$x\ne0$ では
$(U+3)^3=(x^3/y)(U-1)^2$ の根は整で、$U-1$ は単元である。
$\nu(U+3)=\delta/3>\alpha$ より $U/(-3)$ は1に還元される。
$K^2=U$ の符号を $K/(id)$ が1に還元されるように取れば
$\nu(K-id)=\delta/3-\alpha/2$ となる。
また $\nu(H)=\beta/6$ なので証明中の精度25の条件を満たす。
\end{proof}
